\documentclass[10pt,conference,compsocconf]{IEEEtran}

\usepackage{amsmath,amssymb}
\usepackage{booktabs}
\usepackage[hidelinks]{hyperref}

\begin{document}
\title{Linear Models for Heart-Disease Classification}
\author{
  Anthony Glegola \qquad Shu Ito \qquad Yusen Haung\\
  \textit{CS-433 Machine Learning, Project 1}
}
\maketitle

\begin{center}
\footnotesize
\textit{Working draft: updated linear-model results;
final team comparison pending.}
\end{center}

\begin{abstract}
We study classification of self-reported heart-disease history
using the course-provided BRFSS 2015 dataset. We compare least
squares and ridge regression under an earlier partial-feature
baseline and an extended preprocessing pipeline incorporating
categorical encoding, survey-code decoding, and supervised
feature selection. On the same validation split and with fixed
model settings, positive-class F1 increases from 0.3446 to
0.4289 for least squares and from 0.3434 to 0.4296 for ridge.
A subsequent parameter search provides only small additional
gains. We document training-only preprocessing, numerical
rank handling, implementation checks, and reproducible
full-data candidate predictions.
\end{abstract}

\section{Introduction and Data}
The task is to predict the binary \texttt{\_MICHD} label from
health and lifestyle survey features \cite{course,brfss}.
The label describes reported coronary heart disease,
myocardial infarction, or angina history, rather than a
prospectively observed clinical outcome.

The released dataset contains 328,135 labeled observations
and 321 input features, excluding identifiers. Labels are
$y\in\{-1,+1\}$, with approximately 8.83\% positives.
The 109,379 competition test observations have no supplied
labels. This imbalance makes positive-class F1, precision,
and recall important alongside accuracy.

Survey codes require variable-specific interpretation
\cite{codebook}. For example, \texttt{PHYSHLTH}=88 represents
zero days, whereas 77 and 99 represent nonresponse.
Treating these codes as ordinary magnitudes can distort
the feature representation.

\section{Experimental Setup and Preprocessing}
\subsection{Validation protocol}
We shuffle row indices using
\texttt{np.random.default\_rng(42)} and allocate the first
$\lfloor0.25N\rfloor$ observations to validation. This gives
246,102 training and 82,033 validation observations, including
21,737 and 7,238 positives, respectively. The split is not
stratified and is identical across the reported comparisons.

Imputation, scaling, category vocabularies, and supervised
feature selection are fitted only on the training subset.
Validation and competition test observations do not
participate in these fitting steps.

\subsection{Earlier partial-feature baseline}
The initial pipeline uses six columns with one observed
nonmissing value, filling missing entries with $-1$, and
111 ordered or numerical columns. Variable-specific rules
map survey codes to numerical values or missing entries.
Numerical missing values are mean-imputed and standardized
using training statistics.

Including an intercept gives 118 columns. After transformation,
\texttt{CTELNUM1} and \texttt{CELLFON2} are both exactly
$-\texttt{CTELENUM}$. Removing these two redundant columns
gives 116 independent inputs for the linear-model comparison.

\subsection{Extended preprocessing and selection}
We use Shu's updated preprocessing definitions at repository
revision \texttt{f6baef4}. The pipeline includes 5 one-value,
111 ordered, 14 decoded, and 157 categorical raw variables.
It excludes 29 designated variables and 5 decoding candidates
not implemented in that revision.

The numerical groups receive training-mean imputation and
standardization. Decoding converts supported survey
representations, such as encoded weight and frequency
values, into numerical quantities. Categorical variables
are one-hot encoded, including an observed missing-value
category; unseen validation levels produce all-zero entries.
This yields 1,100 columns including the intercept.

Shu's L1-regularized logistic selector uses binary 0/1 labels,
regularization 0.02, step size 0.2, and 2,000 subgradient
iterations. We retain coefficients with absolute value
above 0.001 and always retain the intercept. Fitted on the
training subset, this selects 844 columns.

To support unregularized least squares, we remove numerical
dependencies using training-only pivoted Gram elimination,
with squared relative residual tolerance $10^{-10}$.
Both linear methods then use the same 734 retained columns.
The largest discarded relative residual is
$8.22\times10^{-7}$. Removing columns changes ridge's penalty
representation, so we also evaluate ridge before this step.

\section{Models and Evaluation}
For the transformed matrix $X\in\mathbb{R}^{n\times d}$,
including the intercept, the course loss is
\begin{equation}
L(w)=\frac{1}{2n}\lVert y-Xw\rVert_2^2.
\end{equation}
Setting its gradient to zero gives the least-squares
normal equations. Ridge adds $\lambda\lVert w\rVert_2^2$:
\begin{align}
(X^\top X)w &= X^\top y,\\
(X^\top X+2n\lambda I_d)w &= X^\top y.
\end{align}
We use \texttt{np.linalg.solve} without explicitly computing
an inverse. Both shared functions return \texttt{(w, loss)};
ridge returns the data loss without the penalty
\cite{course}. Following the lab convention, the penalty
includes the intercept.

The regression score $s=Xw$ is not a probability.
We predict $+1$ when $s\geq t$ and $-1$ otherwise.
The fixed comparison uses the earlier threshold $t=-0.55$
and ridge $\lambda=0.001$.

For model selection, we search ridge penalties
$\{10^{-5},10^{-4},10^{-3},10^{-2},10^{-1},1\}$ and thresholds
from $-0.90$ to $0.00$ in steps of $0.01$, maximizing
validation F1:
\begin{equation}
F_1=\frac{2TP}{2TP+FP+FN}.
\end{equation}
Selected thresholds are $-0.57$ for least squares and
$-0.56$ for ridge, with ridge $\lambda=10^{-4}$.
These selected scores reuse validation labels for tuning;
they are not independent test estimates.

\section{Results and Discussion}
\begin{table}[t]
\centering
\caption{Local validation results on the same split.
Partial denotes the earlier representation.
Fixed uses $t=-0.55$ and ridge $\lambda=0.001$.
Tuned uses validation-selected settings.}
\label{tab:results}
\setlength{\tabcolsep}{3pt}
\begin{tabular}{@{}lrrrr@{}}
\toprule
Method & Acc. & Prec. & Recall & F1 \\
\midrule
Always negative & .9118 & .0000 & .0000 & .0000 \\
LS, partial & .8771 & .3255 & .3661 & .3446 \\
Ridge, partial & .8781 & .3272 & .3613 & .3434 \\
\midrule
LS, fixed & .8758 & .3608 & .5287 & .4289 \\
Ridge, fixed & .8752 & .3599 & .5327 & .4296 \\
LS, tuned & .8693 & .3490 & .5569 & .4291 \\
Ridge, tuned & .8728 & .3555 & .5434 & .4298 \\
\bottomrule
\end{tabular}
\end{table}

At fixed model settings, the extended pipeline increases
F1 by 0.0843 for least squares and 0.0862 for ridge
(Table~\ref{tab:results}). Recall increases substantially,
while accuracy decreases slightly. The comparison combines
categorical encoding, decoding, feature selection, and rank
handling; it does not identify the isolated contribution
of any individual transformation.

Parameter search adds only approximately 0.0002 F1 over
the corresponding fixed-setting results. The tuned ridge
and least-squares scores differ by less than 0.001, which
does not establish general superiority on one split.

For tuned least squares, training and validation half-MSE
are 0.129122 and 0.130538. For tuned ridge, they are
0.129242 and 0.130442. These modest gaps do not suggest
severe training-set overfitting, although repeated use of
the validation set limits generalization claims.

Before rank reduction, selected-input ridge at the fixed
settings achieves F1 0.429598, compared with 0.429614 after
reduction. This small observed difference is consistent
with rank handling primarily enabling a stable shared
comparison in this experiment.

The always-negative predictor achieves 91.18\% accuracy
but zero recall and F1. Consequently, accuracy alone would
favor a classifier that detects no positive observations.
Our results highlight the value of reporting several
metrics under class imbalance.

\section{Verification and Reproducibility}
The experiments import the shared methods from
\texttt{implementations.py}. All six required methods are
integrated locally. In the supplied grading environment
(Python 3.9.25, NumPy 1.23.1), 27 local public checks pass,
including the numerical tests and formatting check.
The commit-URL check requires a published final revision.
Four additional tests verify preprocessing consistency,
compact versus dense computations, fitted-state reuse,
and rank handling.

The dataset experiments use Python 3.14.6 and NumPy 2.4.6.
The review notebook is
\texttt{experiments/Yusen\_final\_preprocessing.ipynb};
JSON records preserve settings, metrics, feature counts,
source revisions, and data hashes.

After model selection, we refit preprocessing, L1 selection,
and both linear models on all 328,135 labeled observations.
This produces 1,104 expanded columns, 838 selected columns,
and 731 independent model inputs. These counts differ from
the validation fit because the fitted data and vocabularies
differ.

The command
\texttt{python run.py --model both --predict-only}
reproduces both candidates from included fitted parameters.
Each CSV contains 109,379 predictions. We verify headers,
test-ID order, row counts, and labels, and repeated prediction
produces identical file hashes. These candidates have not
been scored on AIcrowd by this workflow.

\noindent\textbf{Pending team integration.}
Before submission, add the confirmed GD/SGD and logistic
experiment details, identify the team's selected final
pipeline and competition score, and verify its immutable
repository URL \cite{repository}. The comparisons above
cover the linear-model contribution.

\section{Conclusion}
The extended preprocessing pipeline improves positive-class
F1 from approximately 0.34 to 0.43 under matched fixed model
settings. Subsequent parameter tuning provides little
additional improvement. Least squares and ridge perform
similarly, while majority-class accuracy obscures failure
to detect positives. Training-only transformations, explicit
rank handling, and reproducible candidate predictions
provide a documented basis for final team integration.

\begin{thebibliography}{4}
\bibitem{course}
EPFL CS-433 Teaching Team,
``Machine Learning Course: Class Project 1,''
course-provided project specification, 2026.
\bibitem{brfss}
Centers for Disease Control and Prevention,
``2015 BRFSS Survey Data and Documentation.''
[Online]. Available:
\url{https://www.cdc.gov/brfss/annual_data/annual_2015.html}
\bibitem{codebook}
Centers for Disease Control and Prevention,
``BRFSS 2015 Codebook Report:
Land-Line and Cell-Phone Data,'' Aug. 23, 2016.
\bibitem{repository}
A. Glegola, S. Ito, and Y. Haung,
``CS-433 Project 1: Team Repository,'' 2026.
[Online]. Available:
\url{https://github.com/shuito1127/cs433-project1}
\end{thebibliography}
\end{document}
